\documentclass[11pt]{article}
\usepackage[margin=1in]{geometry}
\usepackage{amsmath,amssymb,amsthm,mathtools}
\usepackage{graphicx,xcolor,enumitem,microtype,hyperref}
\hypersetup{colorlinks=true,linkcolor=blue!55!black,citecolor=blue!55!black,urlcolor=blue!55!black}
\newtheorem{theorem}{Theorem}
\newtheorem{lemma}{Lemma}
\newcommand{\W}{\mathcal W}
\newcommand{\E}{\mathcal E}
\newcommand{\Ric}{\operatorname{Ric}}
\newcommand{\Var}{\operatorname{Var}}
\title{\vspace{-2em}A Two-State Counterexample to\\Erbar--Fathi Conjecture 6.9}
\author{Counterexample packet for arXiv:1612.00514}
\date{17 August 2026}

\begin{document}
\maketitle
\begin{center}
\fcolorbox{red!65!black}{red!3}{\parbox{0.91\textwidth}{
\textbf{Status: candidate full counterexample, likely valid.}
The intended assertion that exponential concentration with rate $\rho$
and nonnegative entropic Ricci curvature force a spectral gap
$\lambda\geq C(M)\rho^2$ fails on imbalanced two-state chains. The
dimensionally inconsistent $\rho^{-2}$ exponent printed in the source is
separately false by time rescaling.
}}
\end{center}

\section{Source and exact conjecture}

The source is M.~Erbar and M.~Fathi,
\emph{Poincar\'e, modified logarithmic Sobolev and isoperimetric inequalities
for Markov chains with non-negative Ricci curvature},
arXiv:1612.00514, later JFA 274 (2018), 3056--3089 \cite{EF}.
For a reversible continuous-time chain $(\mathcal X,Q,\pi)$, the source writes
$P(\lambda)$ for
\[
       \Var_\pi f\leq\lambda^{-1}\E(f,f).
\]
Its entropic transport distance $\W$ induces
$d_{\W}(x,y)=\W(\delta_x,\delta_y)$.
Conjecture~6.9 on source page~23 assumes $\Ric(\mathcal X,Q,\pi)\geq0$ and
concentration
\[
 \pi(A_r^c)\leq M e^{-\rho r}
 \quad\text{for every }A\text{ with }\pi(A)\geq\tfrac12,
\]
and asks for a Poincar\'e bound depending only on $M$ and $\rho$.

\begin{figure}[ht]
\centering
\includegraphics[width=0.98\textwidth]{figures/open_problem_crop.png}
\caption{Conjectures 6.9--6.10 on page 23 of the source.}
\end{figure}

The displayed formula in \cite{EF} is $P(C(M)\rho^{-2})$. Since the
Poincar\'e parameter has units of inverse distance squared, the intended
formula is $P(C(M)\rho^2)$. This is also how M\"unch explicitly restates
Conjecture~6.9, as $\lambda\geq C(M)\beta^2$, in his later
Ollivier-curvature analogue \cite[\S6.1]{M}. We disprove the intended
formula nontrivially and the literal printed formula by an independent
scaling argument.

\section{Proof intuition}

Put almost all stationary mass at one endpoint of a two-state chain. Below
the single nonzero metric scale, the concentration tail is just the mass
$p$ of the rare endpoint. It can therefore be written as
$e^{-\rho r}$ with $\rho D=\log(1/p)$.

The intrinsic entropic distance does react to the small atom, but only as
$D^2=O(\log(1/p))$. Consequently $\rho^2$ grows like
$\log(1/p)$, while the generator is normalized to have spectral gap one.
An exact one-dimensional calculation also shows strictly positive entropic
Ricci curvature. Thus atomic concentration at sub-mesh scales is too strong
to quantitatively control the spectral gap.

\section{The two-state geometry}

Fix $0<p\leq1/4$, put $q=1-p$, and consider
\[
 \mathcal X=\{0,1\},\qquad
 Q(0,1)=p,\quad Q(1,0)=q,\qquad
 \pi(0)=q,\quad\pi(1)=p.
\]
Detailed balance holds, and the nonzero eigenvalue of $-L$ is $p+q=1$.
Hence the optimal Poincar\'e parameter is
\begin{equation}\label{eq:gap}
                         \lambda_p=1.
\end{equation}

Parameterize a probability measure by $\mu_m=(1-m,m)$, $0\leq m\leq1$.
Its density relative to $\pi$ is
\[
 \rho_m(0)=\frac{1-m}{q},\qquad \rho_m(1)=\frac m p.
\]
Let $\theta(a,b)=(a-b)/(\log a-\log b)$ be the logarithmic mean and put
\[
 \ell_p(m)=\log\frac{mq}{p(1-m)},\qquad
 K_p(m)=pq\,\theta\!\left(\frac{1-m}{q},\frac m p\right)
        =\frac{m-p}{\ell_p(m)}.
\]
The two directed edge terms and the factor $1/2$ in the action give
$\mathcal A(\rho_m,\psi)=K_p(m)(\psi(1)-\psi(0))^2$.
The continuity equation gives
$\dot m=K_p(m)(\psi(1)-\psi(0))$ up to an irrelevant sign. Therefore the
Riemannian line element and endpoint distance are
\begin{equation}\label{eq:metric}
 ds^2=\frac{dm^2}{K_p(m)},\qquad
 D_p=d_{\W}(0,1)
 =\int_0^1\sqrt{\frac{\ell_p(m)}{m-p}}\,dm.
\end{equation}
The quotient is positive, extends continuously at $m=p$, and has integrable
endpoint singularities.

\begin{lemma}[Positive entropic curvature]\label{lem:ric}
The chain satisfies $\Ric(\mathcal X,Q,\pi)\geq1/2$.
\end{lemma}

\begin{proof}
The entropy along the line is
\[
 H_p(m)=m\log\frac mp+(1-m)\log\frac{1-m}{q},
 \quad H_p'(m)=\ell_p(m),\quad
 H_p''(m)=\frac1{m(1-m)}.
\]
For the one-dimensional metric tensor $G=1/K_p$, the Riemannian Hessian
relative to the metric equals
\[
 \frac{\operatorname{Hess}_G H_p}{G}
 =K_p\left(H_p''-\frac{G'}{2G}H_p'\right)
 =\frac12\left(1+\frac{K_p(m)}{m(1-m)}\right)\geq\frac12.
\]
This is precisely $1/2$-convexity of entropy along $\W$-geodesics in the
interior. Continuity of the entropy and finiteness of \eqref{eq:metric}
extend it to endpoint measures.
\end{proof}

\begin{lemma}[Logarithmic diameter bound]\label{lem:D}
For $0<p\leq1/4$ and $L=\log(1/p)$,
\[
                            D_p\leq7\sqrt L,
 \qquad\text{hence}\qquad D_p^2\leq49L.
\]
\end{lemma}

\begin{proof}
Write $g(u)=\log(u/(1-u))$, so the integrand in \eqref{eq:metric} is
$\sqrt{(g(m)-g(p))/(m-p)}$. Split the integral at $p,2p,1/2$.
On $(0,p)$, since $(1-s)^{-1}\leq4/3$ for $0<s\leq p$,
\[
 \int_0^p\sqrt{\frac{g(p)-g(m)}{p-m}}\,dm
 \leq\sqrt{\frac{4p}{3}}
 \int_0^1\sqrt{\frac{\log(1/u)}{1-u}}\,du
 <\frac4{\sqrt3}\sqrt p.
\]
The last universal integral is less than $2$ (split at $u=1/2$).
On $(p,2p)$, the mean-value formula and
$g'(u)=1/[u(1-u)]$ give an integrand at most $\sqrt{2/p}$, so this piece is
at most $\sqrt{2p}$.

For $2p\leq m\leq1/2$,
$g(m)-g(p)\leq L$ and $m-p\geq m/2$, whence this piece is at most
$2\sqrt L$. Finally, for $1/2\leq m<1$, we have $m-p\geq1/4$ and
\[
 g(m)-g(p)\leq L+\log\frac1{1-m}.
\]
Using $\sqrt{a+b}\leq\sqrt a+\sqrt b$, this last piece is at most
\[
 \sqrt L+2\int_0^{1/2}\sqrt{\log(1/u)}\,du
 \leq\sqrt L+\sqrt\pi.
\]
Since $L\geq\log4$, the sum of the four bounds is less than
$7\sqrt L$.
\end{proof}

\section{Counterexample to the intended conjecture}

\begin{theorem}\label{thm:counter}
There is no positive constant $C(1)$ such that every finite reversible chain
with nonnegative entropic Ricci curvature and exponential concentration
profile $e^{-\rho r}$ satisfies $P(C(1)\rho^2)$.
In particular, the intended form of Conjecture~6.9 is false.
\end{theorem}

\begin{proof}
For the two-state chain, every set of $\pi$-mass at least $1/2$ is either
$\{0\}$ or the whole space. For $A=\{0\}$,
\[
 \pi(A_r^c)=
 \begin{cases}
 p,&0\leq r<D_p,\\
 0,&r\geq D_p.
 \end{cases}
\]
Thus, with
\[
                  \rho_p=\frac{\log(1/p)}{D_p},
\]
the measure has concentration profile $e^{-\rho_p r}$ (with $M=1$).
Lemmas~\ref{lem:ric} and \ref{lem:D} give the curvature hypothesis and
\[
                     \rho_p^2
     =\frac{L^2}{D_p^2}\geq\frac{L}{49}\longrightarrow\infty.
\]
But \eqref{eq:gap} says that $P(\lambda)$ holds only for $\lambda\leq1$.
The conjectured $P(C(1)\rho_p^2)$ is therefore impossible for all small $p$.
\end{proof}

\section{The literal printed exponent}

\begin{theorem}\label{thm:literal}
The formula $P(C(M)\rho^{-2})$ printed in Conjecture~6.9 is also false for
every positive $C(1)$.
\end{theorem}

\begin{proof}
Fix one $p\in(0,1/2)$ and multiply both transition rates by $s>0$.
The curvature remains nonnegative, the spectral gap becomes
$\lambda_s=s$, and the entropic action is multiplied by $s$, so
$D_s=D_p/\sqrt s$. The same exact concentration argument gives
\[
 \rho_s=\frac{\log(1/p)}{D_s}
       =\sqrt s\,\frac{\log(1/p)}{D_p}.
\]
The literal assertion would require
\[
 s=\lambda_s\geq C(1)\rho_s^{-2}
   =\frac{C(1)D_p^2}{s\log^2(1/p)},
\]
which fails as $s\downarrow0$.
\end{proof}

\section{Verification, novelty, and limitations}

\paragraph{Verification.}
The proof uses only the exact two-state action and elementary integral
bounds. The included regression script independently evaluates $D_p$ by
quadrature, checks the pointwise curvature formula on dense grids, verifies
$D_p^2<49\log(1/p)$ down to $p=10^{-8}$, and checks the rescaling law.
It is not used as proof. The strict definition
$A_r^c=\{d(x,A)>r\}$ causes no endpoint issue: the tail is zero at $r=D_p$.

\paragraph{Bounded novelty search.}
The four run indexes were searched for arXiv:1612.00514, Conjecture~6.9,
entropic curvature, exponential concentration, spectral gap, and two-state
chains. Exact-phrase, title/citation, author, two-point, and close-variant
web/arXiv searches through 17 August 2026 found no counterexample to the
original entropic-curvature statement. M\"unch \cite{M} explicitly claims
affirmative analogues under nonnegative Ollivier curvature, but his results
use a different curvature/metric framework and discrete-scale estimates;
they do not establish the source's entropic-$d_{\W}$ conjecture. Novelty is
plausible, not certified.

\paragraph{Scope and review.}
This completely disproves Conjecture~6.9 under either exponent reading. It
does not disprove or prove Conjecture~6.10; the same two-state family does not
produce an MLSI contradiction. The mechanism exploits an atom below the
first nonzero metric scale, suggesting that a repaired theorem must truncate
the concentration rate at that scale or impose a non-atomic/geodesic
regularity condition. Human review should focus on the normalization in
\eqref{eq:metric}, the one-dimensional Hessian calculation, and the source
exponent interpretation.

\begingroup
\sloppy
\raggedright
\begin{thebibliography}{9}
\bibitem{EF}
M.~Erbar and M.~Fathi,
\emph{Poincar\'e, modified logarithmic Sobolev and isoperimetric inequalities
for Markov chains with non-negative Ricci curvature},
J. Funct. Anal. \textbf{274} (2018), 3056--3089;
\href{https://arxiv.org/abs/1612.00514}{arXiv:1612.00514}.

\bibitem{M}
F.~M\"unch,
\emph{Ollivier curvature, isoperimetry, concentration, and log-Sobolev
inequality},
\href{https://arxiv.org/abs/2309.06493}{arXiv:2309.06493}, 2023.
\end{thebibliography}
\endgroup
\end{document}
